% 原书 p.12–p.15（PDF p.20–p.23）
% >>> p020
\section{辛同胚（Symplectomorphism）}

\begin{keytakeaways}
\begin{itemize}
  \item 几何图像：辛同胚是保持辛形式的微分同胚，它是「辛流形之间的同构」；
        $\mathrm{Symp}(M,\omega)$ 是 $\Diff(M)$ 的子群，其单位连通分支记作
        $\mathrm{Symp}_0(M,\omega)$。
  \item 数学工具：线性化——把 $\mathrm{Symp}$ 的问题化为矩阵群 $\mathrm{Sp}(2n)$
        （条件 (1.11)）；极分解 $A=PO$ 与特征多项式恒等式 (1.12)。
  \item 本节结论：$\mathrm{Sp}(2n)$ 同伦等价于 $U(n)$，故 $\pi_1(\mathrm{Sp}(2n))=\Z$；
        辛群与保体积群是两类不同的群（Gromov 非压缩定理 1.5.1 给出首个显式区别）。
  \item 定量图像：辛球不能以小于其「半径」的辛柱压缩——这是辛刚性最著名的表现，
        也是「辛」与「保体积」分野的定量标志。
\end{itemize}
\end{keytakeaways}

与微分流形之间的微分同胚类似，辛同胚把辛流形彼此连接起来，并帮助我们对辛结构进行分类。

\begin{wdefn}{1.5.1}
微分同胚 $f:(M,\omega_M)\to(N,\omega_N)$ 称为\textbf{辛同胚}（symplectomorphism），若
$f^*\omega_N=\omega_M$。$(M,\omega)$ 到自身的所有辛同胚构成的集合记作 $\Symp(M,\omega)$。
\end{wdefn}

容易验证 $\Symp(M,\omega)$ 是一个群，并且是 $M$ 上微分同胚群（通常记作 $\Diff(M)$）的子群。
注意 $\Symp(M,\omega)$ 可能含有多个连通分支，而 $\Symp(M,\omega)$ 的单位连通分支记作
$\Symp_0(M,\omega)$。此外，直接由定义 1.5.1，任何辛同胚 $f\in\Symp(M,\omega)$ 都保持体积形式，
即 $f^*\omega^n=\omega^n$。

\begin{wrem}{1.5.1}
与定义 1.5.1 相比，有时人们关心一个更为灵活的定义：$f^*\omega_N=c\cdot\omega_M$，其中 $c>0$。
由此得到的研究对象称为「\textbf{共形}」辛几何（conformal symplectic geometry）。
\end{wrem}

群 $\Symp(M,\omega)$ 的性质相当神秘，即使在 $(M,\omega)=(\R^{2n},\omega_{\mathrm{std}})$ 时也是如此。
让我们先看 $\Symp(\R^{2n},\omega_{\mathrm{std}})$ 中那些线性的元素，即由矩阵表示的元素。

\begin{wexam}{1.5.1}
用 $M_{2n\times2n}(\R)$ 记 $\R$ 上 $2n\times2n$ 矩阵构成的群。按定义 1.5.1，
$A\in M_{2n\times2n}(\R)$ 是线性辛同胚，若 $A^*\omega_{\mathrm{std}}=\omega_{\mathrm{std}}$。
在 $\R^{2n}$ 的有序基 $(x_1,\cdots,x_n,y_1,\cdots,y_n)$ 之下，这一关系等价于
\begin{equation}
  A^{\mathrm T}\Omega A=\Omega
  \qquad\text{其中}\qquad
  \Omega=\begin{pmatrix}0&\mathbf{1}_{n\times n}\\[1.5mm]
    -\mathbf{1}_{n\times n}&0\end{pmatrix}.
  \tag{1.11}
\end{equation}
$(\R^{2n},\omega_{\mathrm{std}})$ 上所有线性辛同胚构成的集合记作 $\mathrm{Sp}(2n)$，它是一个群。
注意有些辛几何学家只与 $\mathrm{Sp}(2n)$ 打交道。\hfill$//$
\end{wexam}

一般地，条件 (1.11) 并不总是容易验证，尤其是矩阵的规模很大时。有关更有效的判别准则，
见 [153] 中的练习 1.1.13。
% [Tutor Note: 记号提醒——(1.11) 中的 $\mathbf{1}_{n\times n}$ 是 $n$ 阶单位矩阵，而
% $\Omega$ 是 $\omega_{\mathrm{std}}$ 在有序基 $(x_1,\cdots,x_n,y_1,\cdots,y_n)$ 下的矩阵：
% $\omega_{\mathrm{std}}=\sum_i\dd x_i\wedge\dd y_i$ 给出 $\Omega=\begin{pmatrix}0&I\\-I&0\end{psmallmatrix}$。
% 若把基的顺序改成 $(x_1,y_1,\cdots,x_n,y_n)$，则 $\Omega$ 变成分块对角、每块为
% $\begin{pmatrix}0&1\\-1&0\end{psmallmatrix}$，两式只差一个置换矩阵的共轭；对照文献时须先确认基序。
% 由 $A^{\mathrm T}\Omega A=\Omega$ 立得 $\det A=\pm1$，再结合保体积性（见下文 (i)）才得 $\det A=1$。]

% >>> p021
\begin{wque}{1.5.1}
关于 $\mathrm{Sp}(2n)$ 中的元素 $A$，我们能知道多少？
\end{wque}

下面列出一些初等性质。

\begin{enumerate}[label=(\roman*)]
  \item $A$ 的行列式为 $1$。注意由条件 (1.11)，我们只能推出
        $\det(A^{\mathrm T}\Omega A)=\det(\Omega)$，这蕴涵 $\det(A)^2=1$。要证明 $\det(A)=1$
        确实成立，有各种办法。一个快速的办法是利用 $A$ 的保体积性质：由于
        $A^*\omega_{\mathrm{std}}^n=\omega_{\mathrm{std}}^n$，我们有 $\det(A)=1$。我们留有兴趣的
        读者去寻找尽可能多的途径来证明这一点。
  \item $A$ 的特征多项式，记作 $P_A(\lambda)$，满足
    \begin{equation}
      P_A(\lambda)=\lambda^{2n}P_A(\lambda^{-1}).
      \tag{1.12}
    \end{equation}
    因此，$A$ 的任意特征值 $\lambda\in\C$ 总是成组出现：$\lambda,\lambda^{-1},\bar\lambda,1/\bar\lambda$。
    注意 $\bar\lambda$ 之所以出现，是因为 $P_A(\lambda)$ 是实系数多项式，故任何复根都伴随其共轭
    一同出现。下面我们证明 (1.12)。
\[
  \begin{aligned}
    P_A(\lambda)&=\det(A-\lambda\mathbf{1})
      =\det\big(-\Omega(A^{\mathrm T})^{-1}\Omega+\lambda\Omega\cdot\Omega\big)\\
      &=\det(\Omega)\cdot\det(\Omega)\cdot\det\big(-(A^{\mathrm T})^{-1}+\lambda\mathbf{1}\big)\\
      &=\det\big(-(A^{\mathrm T})^{-1}+\lambda\mathbf{1}\big)\\
      &=\det\big((-A^{\mathrm T})^{-1}(\mathbf{1}-\lambda A^{\mathrm T})\big)\\
      &=\det(\mathbf{1}-\lambda A)\\
      &=(-\lambda)^{2n}\det(A-\lambda^{-1}\mathbf{1})\\
      &=\lambda^{2n}\det(A-\lambda^{-1}\mathbf{1})=\lambda^{2n}P_A(\lambda^{-1}).
  \end{aligned}
\]
\end{enumerate}
% [Tutor Note: 第二行到第三行用到 $\det(\Omega)=\pm1$ 的平方必为 1——按 (1.11) 的约定
% $\det\Omega=1$（$\Omega$ 是置换矩阵与 $\mathrm{diag}(J,\cdots,J)$ 的乘积，$J=\begin{psmallmatrix}0&1\\-1&0\end{psmallmatrix}$，
% 每个 $J$ 的行列式为 1），故 $\det(\Omega)^2=1$。第四行用 $-(A^{\mathrm T})^{-1}+\lambda\mathbf{1}
% =(-A^{\mathrm T})^{-1}(\mathbf{1}-\lambda A^{\mathrm T})$，其行列式乘以
% $\det((-A^{\mathrm T})^{-1})=\det(-(A^{\mathrm T})^{-1})=(-1)^{2n}\det(A)^{-1}=1$。最后一步
% $\det(\mathbf{1}-\lambda A)=(-\lambda)^{2n}\det(A-\lambda^{-1}\mathbf{1})$ 是因为
% $\mathbf{1}-\lambda A=-\lambda(A-\lambda^{-1}\mathbf{1})$，而维数 $2n$ 为偶数使 $(-\lambda)^{2n}=\lambda^{2n}$。
% 这一结构正是「辛特征值四元组 $\{\lambda,\lambda^{-1},\bar\lambda,\bar\lambda^{-1}\}$」的来源。]

\begin{wexer}{1.5.1}
求出所有满足 $A\in\mathrm{Sp}(2)$ 的 $A$。对 $\mathrm{Sp}(4)$ 作同样的事情。
\end{wexer}
% [Tutor Note: 原句「Find all $A$ such that $A\in\mathrm{Sp}(2)$」表述上略嫌循环，实意是
% 「确定 $\mathrm{Sp}(2)$ 的全部元素」，答案为所有行列式为 1 的 $2\times2$ 实矩阵，即
% $\mathrm{Sp}(2)=\mathrm{SL}(2,\R)$（低维的例外同构）；而 $\mathrm{Sp}(4)$ 的分类要复杂得多，
% 需用到 (ii) 中特征值的四元组结构逐一讨论。]

\begin{enumerate}[label=(\roman*),start=3]
  \item 回忆对任意可逆的 $A\in M_{2n\times2n}(\R)$，其极分解为
\[
  A=\sqrt{AA^{\mathrm T}}\cdot\left(\sqrt{AA^{\mathrm T}}\right)^{-1}A,
\]
其中 $P:=\sqrt{AA^{\mathrm T}}$ 是正定对称矩阵，而 $O:=\left(\sqrt{AA^{\mathrm T}}\right)^{-1}A$
是正交矩阵。下面这个事实有用。
\end{enumerate}

\whead{事实}{1.5.1}\quad 若 $A\in\mathrm{Sp}(2n)$，则上面的 $P$ 与 $O$ 也都是线性辛同胚。

事实上，只需证明 $P\in\mathrm{Sp}(2n)$，因为 $\mathrm{Sp}(2n)$ 是群。此外容易验证：若
$P\in\mathrm{Sp}(2n)$，则对任意 $\alpha>0$ 有 $P^\alpha\in\mathrm{Sp}(2n)$。这里 $P^\alpha$
是通过其特征值的幂来定义的（由于 $P$ 实对称，它完全由其特征值所决定）。这些合起来蕴涵：
任何线性辛同胚 $A\in\mathrm{Sp}(2n)$ 都可以在 $\mathrm{Sp}(2n)$ 内部如下形变，
\[
  A=PO\ \xrightarrow{\ \text{沿 }P^\alpha O\text{ 形变}\ }\ O\in\mathrm{Sp}(2n)\cap O(2n).
\]
% [Tutor Note: 这里省略了 $P\in\mathrm{Sp}(2n)$ 的验证：由 $A^{\mathrm T}\Omega A=\Omega$ 与
% $AA^{\mathrm T}=A^{\mathrm T}A$（二者都等于 $-A^2$）可得 $P=\sqrt{AA^{\mathrm T}}$ 也满足
% $P^{\mathrm T}\Omega P=\Omega$；进而 $P$ 的特征值全为正实数且以倒数对出现，故 $P^\alpha$ 仍辛。
% 一旦 $P,O$ 都辛，$P^\alpha O$ 便给出 $\mathrm{Sp}(2n)$ 内连接 $A=PO$ 与 $O$ 的道路，
% 这就是「$\mathrm{Sp}(2n)$ 形变收缩到 $\mathrm{Sp}(2n)\cap O(2n)$」的实质，也是下文
% $\pi_k(\mathrm{Sp}(2n))=\pi_k(U(n))$ 的来源。]

% >>> p022
换言之，$\mathrm{Sp}(2n)$ 同伦等价于 $\mathrm{Sp}(2n)\cap O(2n)$。同时，$\mathrm{Sp}(2n)\cap O(2n)$
中的任何元素都具有如下特殊形式：
\[
  \begin{pmatrix}X&-Y\\ Y&X\end{pmatrix}
  \qquad\text{其中}\qquad
  X^{\mathrm T}Y=Y^{\mathrm T}X
  \quad\text{且}\quad
  X^{\mathrm T}X+Y^{\mathrm T}Y=\mathbf{1}.
\]
这可以与矩阵 $X+\sqrt{-1}\,Y\in U(n)$ 等同起来。总之，我们得到结论：
$\pi_k(\mathrm{Sp}(2n))=\pi_k(U(n))$ 对任意 $k\in\N$ 成立。特别地，
$\pi_1(\mathrm{Sp}(2n))=\Z$。

我们已经看到，群 $\Symp(M,\omega)$ 包含于 $\Diff_{\mathrm{vol}}(M)$，即 $M$ 的所有保体积微分同胚
构成的群。下面这个由 Gromov [99] 得到的定理 1.5.1 是现代辛几何的里程碑。它通过发现
$\Diff_{\mathrm{vol}}(M)\setminus\Symp(M,\omega)$ 中存在元素，把这两个群区分开来。因此，辛几何
（辛拓扑）成为一个相对独立于经典微分拓扑的自足学科。

让我们引入一些记号：
\[
  B^{2n}(R):=\left\{(x_1,y_1,\cdots,x_n,y_n)\in\R^{2n}
  \ \Big|\ \sum_{i=1}^{n}(x_i^2+y_i^2)\le R\right\}
\]
以及
\begin{equation}
  Z^{2n}(r):=B^2(r)\times\R^{2n-2}
  \qquad\text{其中}\qquad
  B^2(r)=\{(x_1,y_1)\in\R^2\mid x_1^2+y_1^2\le r\}.
  \tag{1.13}
\end{equation}
$B^{2n}(R)$ 与 $Z^{2n}(r)$ 都是 $\R^{2n}$ 中的闭区域。前者称为\textbf{辛球}（symplectic ball），
后者称为\textbf{辛柱}（symplectic cylinder）。

\begin{wthm}{1.5.1}（Gromov 非压缩定理）
辛球 $B^{2n}(R)$ 经由辛同胚嵌入辛柱 $Z^{2n}(r)$，当且仅当 $R\le r$。
\end{wthm}
% [Tutor Note: 特别注意 (1.13) 的常数约定：$B^{2n}(R)$ 是「半径平方」为 $R$ 的球（即半径为
% $\sqrt R$），而 $Z^{2n}(r)$ 的第二因子取 $B^2(r)$（半径 $\sqrt r$），定理中的比较是
% $R\le r$ 而非 $\sqrt R\le r$。许多文献把球写成 $B^{2n}(r)$（半径为 $r$，即
% $\sum(x_i^2+y_i^2)\le r^2$），此时非压缩定理应读作 $r\le R$；对照文献时务必核对这一约定，
% 否则结论的方向会颠倒。]

\begin{wexer}{1.5.2}
若把 $Z^{2n}(r)$ 的定义改为
\[
  \{(x_1,x_2)\in\R^2\mid x_1^2+x_2^2\le r\}\times\R^{2n-2},
\]
则证明：对任意 $r>0$，$B^{2n}(R)$ 都能辛嵌入 $Z^{2n}(r)$。
\end{wexer}
% [Tutor Note: 本题的要点是：非压缩定理之所以成立，是因为辛柱中被约束的两个坐标构成一个
% \emph{复}方向（$x_1,y_1$ 是一对共轭坐标）；若换成两个同类的实坐标 $(x_1,x_2)$，则可用
% 一个挤压缩放（沿 $x_2$ 方向压缩、同时沿 $y_2$ 方向拉伸，保持 $\dd x_i\wedge\dd y_i$ 之和不变）
% 把整个球装进任意细的柱中。这正是定理 1.5.1 的「复本质」。]

这个练习表明定理 1.5.1 在其本质上是复的：$Z^{2n}(r)$ 中起约束作用的坐标必须构成一个复平面！
事实上，定理 1.5.1 的证明深刻地依赖于对 $J$-全纯曲线的研究。更多细节见 [154]。如今已有
多种证明它的办法。借助更高级的机器，某些证明既快又干净（见 6.4 节）。

让我们用另外两个问题来结束本节。

\begin{wque}{1.5.2}
\begin{enumerate}[label=(\roman*)]
  \item 如何判别一个微分同胚是辛同胚？
  \item 如何构造辛同胚？
\end{enumerate}
\end{wque}

% >>> p023
我们将在接下来的两节中回答这两个问题。对 (i) 的一个快速回答是：人们可以把「是辛同胚」这一
性质转化为子流形的一个几何性质：是拉格朗日的（见 1.6 节）；对 (ii)，人们诉诸哈密顿动力学，
所得的流帮助我们构造出一族特殊的辛同胚（见 1.7 节）。
